% Sampo notes
% Tämä on mun gradu. Tässä alussa vähän löpinää siitä, mistä kyse ja varmaan jotain linkkejä

%%% LAITOKSEN VANHAT GRADUOHJEET
% https://blogs.helsinki.fi/geotiede/files/2008/03/Pro-gradu-ohjeet.pdf
% FLAMMAN GRADUOHJEET:
%%% https://flamma.helsinki.fi/portal/units/matlu?_nfpb=true&_pageLabel=P9806621561360920311893&contentId=HY358060&placeId=HY357991&lang=fi
% Multifile LaTeX project:
%%% https://www.overleaf.com/learn/latex/Multi-file_LaTeX_projects
% How to write a thesis in LaTeX
%%% https://www.overleaf.com/learn/latex/How_to_Write_a_Thesis_in_LaTeX_(Part_1):_Basic_Structure

\documentclass[a4paper,11pt]{article}

% Set margins
\usepackage{geometry}
\geometry{
    a4paper,
    total={170mm,257mm},
    left=20mm,
    right=20mm,
    }

% Font and layout settings
\renewcommand{\baselinestretch}{1.5}        % set line spacing
\usepackage[finnish,english]{babel}         % set languages. Last in list is main language
\usepackage[T1]{fontenc}
\usepackage[utf8]{inputenc}
\usepackage[font=small,labelfont=bf]{caption}
\setlength{\abovecaptionskip}{15pt plus 3pt minus 2pt}
\setlength{\belowcaptionskip}{15pt plus 3pt minus 2pt}
\usepackage[document]{ragged2e}             % package enables justifying of text. \justify
\usepackage{lastpage}                       % get amount of pages

% Datetime settings
\usepackage[ddmmyyyy]{datetime}             % change format of datetime
%\renewcommand{\dateseparator}{.}
\newdateformat{mydate}{%
\onedigit{\THEDAY}.\onedigit{\THEMONTH}.\THEYEAR} % remove unnecessary zeroes in date
\newdateformat{monthyeardate}{%
    \monthname[\THEMONTH] \THEYEAR}         % mm/yyyy date format for English abstract
\newcommand*{\printfindate}{%               % Get Finnish months and year for Finnish abstract
   \ifcase \month\or Tammikuu\or Helmikuu\or Maaliskuu\or Huhtikuu\or Toukokuu\or Kesäkuu\or Heinäkuu\or Elokuu\or Syyskuu\or Lokakuu\or Marraskuu\or Joulukuu\fi \space \number\year}

% location of files
\usepackage{graphicx}
\graphicspath{ {Images/} }

% references and reference styles
\usepackage{enumitem}% http://ctan.org/pkg/enumitem
\usepackage{natbib}
\bibliographystyle{agsm}
\setcitestyle{authoryear,open={},close={}}  % Removes citation parentheses

% other packages
\usepackage{alphabeta}
\usepackage{amsmath} 
\usepackage{subfigure}
\usepackage{gensymb}
\usepackage[colorinlistoftodos]{todonotes}
\usepackage{hyperref} 
\usepackage{amsthm} 
\usepackage{float}
\usepackage{pdfpages}
\usepackage{minted}
\usepackage{enumitem}
\usepackage{listings}
\usepackage{color}
\usepackage{import}

% tables
\usepackage[showframe]{geometry}
\usepackage{tabularx}
\usepackage{makecell}                       % enable thicker vline
\usepackage{cellspace}
\setlength{\cellspacetoplimit}{6pt}
\setlength{\cellspacebottomlimit}{6pt}
\addparagraphcolumntypes{X}
\newcolumntype{?}{!{\vrule width 1.5pt}}    % thicker vline for tabularx
\newcommand\Tstrut{\rule{0pt}{2.2ex}}       % "top" strut
\newcommand\Bstrut{\rule[-4.0ex]{0pt}{0pt}} % "bottom" strut
\newcommand{\TBstrut}{\Tstrut\Bstrut}       % top & bottom struts

% translator command setup
% this enables us to use multilanguage strings which works with package babel
\newcommand{\newtranslatorcommand}[3]{%
    \newcommand{#1}{%
        \iflanguage{english}{#2}{%
        \iflanguage{finnish}{#3}{%
        }}%
    }%
}

% Title page and abstract settings
\usepackage{setspace}                       % different line spacing than document default
\newcommand{\myname}{                      % THIS STRING REPRESENTS THESIS AUTHOR IN THE DOCUMENT
    Sampo Vesanen%
}

% THESIS TITLE MULTILANGUAGE STRINGS
% Multilanguage code is adapted from here:
% https://tex.stackexchange.com/questions/42784/language-dependent-custom-command
\newtranslatorcommand{\mytitle}{%
    Parking private cars and accessibility in Helsinki Capital Region%
}{%
    Henkilöauton pysäköinti ja spatiaalinen saavutettavuus pääkaupunkiseudulla%
}

% Choose correct date format for abstract
\newtranslatorcommand{\translatedate}{%
    \monthyeardate\today%
}{%
    \printfindate%
}

% PAGE COUNT MULTILANUAGE STRING
\newtranslatorcommand{\mypages}{%
    \pageref{myLastPage} pages and appendices%
}{%
    \pageref{myLastPage} sivua ja liitteet%
}

% UNIVERSITY MULTILANGUAGE STRING
\newtranslatorcommand{\myuni}{%
    University of Helsinki%
}{%
    Helsingin yliopisto%
}

% FACULTY MULTILANGUAGE STRING
\newtranslatorcommand{\myfaculty}{%
    Faculty of Science%
}{%
    Matemaattisluonnontieteellinen tiedekunta%
}

% DEPARTMENT MULTILANGUAGE STRING
\newtranslatorcommand{\mydept}{%
    Department of Geosciences and Geography%
}{%
    Geotieteiden ja maantieteen osasto%
}

% GEOGRAPHY MULTILANGUAGE STRING
\newtranslatorcommand{\mygeography}{%
    Geography
}{%
    Maantiede
}

% GEOINFORMATICS MULTILANGUAGE STRING
\newtranslatorcommand{\mygis}{%
    Geoinformatics%
}{%
    geoinformatiikka%
}

% THESIS TYPE MULTILANGUAGE STRING
\newtranslatorcommand{\mythesis}{%
    Master's thesis%
}{%
    Pro gradu -tutkielma%
}

% CHOOSE CORRECT ABSTRACT LANGUAGE
\newtranslatorcommand{\myabstract}{%
    \import{sections/}{abstract_en.tex}%
}{%
    \import{sections/}{abstract_fi.tex}%
}

% KEYWORDS MULTILANGUAGE STRING
\newtranslatorcommand{\mykeywords}{%
    Parking, accessibility, GIS, Python%
}{%
    Pysäköinti, saavutettavuus, GIS, Python%
}

% LIBRARY MULTILANGUAGE STRING
\newtranslatorcommand{\mylibrary}{%
    Kumpula science library%
}{%
    Kumpulan tiedekirjasto%
}

% OTHER INFO MULTILANGUAGE STRING
\newtranslatorcommand{\myotherinfo}{%
    -%
}{%
    -%
}

\title{\mytitle}
\author{\myname}
\date{\today}


% BEGIN DOCUMENT PROPER
\begin{document}

% TITLE PAGE
\begin{titlepage}{
    \setstretch{1.0}
    \centering
    \begin{figure}[t]
    \centering
    \includegraphics[scale=0.4]{HY_logo}
    \end{figure}
    
    \mythesis \par
    \mygeography \par
    \mygis \par
    
    \bigskip
    \MakeUppercase{\mytitle}
    
    \bigskip
    \myname
    
    \mydate\today
    
    % Fills empty space between top and bottom of page
    \vfill
    
    Supervisor: \par
    Tuuli Toivonen \par
    \bigskip
    \bigskip
    \MakeUppercase\myuni \par
    \MakeUppercase\myfaculty \par
    \MakeUppercase\mydept \par
    \MakeUppercase{\mygeography} \par
    PL 64 (Gustaf Hällströmin katu 2) \par
    00014 Helsingin yliopisto \par
}
\end{titlepage}

%%% 1ST ABSTRACT PAGE, ENGLISH
\newpage
\import{sections/}{abstract.tex}

%%% 2ND ABSTRACT PAGE, FINNISH
\newpage
\begin{otherlanguage*}{finnish}
\import{sections/}{abstract.tex}
\end{otherlanguage*}

% TABLE OF CONTENTS
\newpage
\selectlanguage{english}
\thispagestyle{empty}
\renewcommand*\contentsname{Table of contents} % Set custom name for ToC
\tableofcontents

% INTRODUCTION AND RESEARCH QUESTIONS
\newpage
\setcounter{page}{1} % set page numbering to start from here
\section{Introduction}
\justify
Tässä kappaleessa ei alaotsikoita.
%\import{/sections}{introduction.tex}

The specific research questions for this thesis are:

\begin{spacing}{1}
\begin{enumerate}[label=\Roman*] % use Roman numerals, uses enumitem package
  \item What are the spatial differences in the time that it takes to find and park one’s car in the study area?
  \item If spatial differences are detected, what explains them?
  \item What is the significance of the process of parking one’s car to the overall travel time?
\end{enumerate}
\end{spacing}

% THEORETICAL BACKGROUND
\newpage
\section{Background}
\subsection{Parking studies}
\subsection{Parking time estimations}

% DATA AND METHODS
\newpage
\section{Data and methods}
\subsection{Parking survey}
\justify
To test if parking really is hard in Helsinki, I created a survey in JavaScript and Leaflet. It was actually the second time I was doing the survey as the first one, done with Survey123 didn't get approved. That really sucked man.

% RESULTS
\newpage
\section{Results}

% DISCUSSION
\newpage
\section{Discussion}

% REFERENCES
% Connect Mendeley references to Overleaf bib file
\newpage
\label{myLastPage} % this label defines the last page to be included in page count
\section{References}
\renewcommand{\refname}{}
\bibliography{references.bib}

% APPENDICES
\newpage
\pagenumbering{gobble} % suppress appendices page numbering from page and table of contents 
\section{Appendices}

\end{document}
